% ১৪.১০ Input Output Statement
\section{Input/Output Statement}

\begin{learningbox}
এই টপিক শেষে শিক্ষার্থী \texttt{printf} ও \texttt{scanf} ব্যবহার করে input-output করতে পারবে।
\end{learningbox}

\begin{definitionbox}{\texttt{printf}}
C language-এ output দেখানোর জন্য \texttt{printf()} function ব্যবহৃত হয়।
\end{definitionbox}

\begin{definitionbox}{\texttt{scanf}}
Keyboard থেকে input নেওয়ার জন্য \texttt{scanf()} function ব্যবহৃত হয়।
\end{definitionbox}

\begin{verbatim}
int n;
scanf("%d", &n);
printf("Number = %d", n);
\end{verbatim}

\begin{center}
\begin{tabular}{|p{0.18\textwidth}|p{0.30\textwidth}|p{0.30\textwidth}|}
\hline
\textbf{Specifier} & \textbf{Data type} & \textbf{Example} \\
\hline
\texttt{\%d} & int & \texttt{10} \\
\hline
\texttt{\%f} & float & \texttt{3.5} \\
\hline
\texttt{\%c} & char & \texttt{A} \\
\hline
\texttt{\%s} & string & \texttt{ICT} \\
\hline
\end{tabular}
\end{center}

\begin{mistakebox}{সাধারণ ভুল}
\texttt{scanf}-এ variable-এর আগে সাধারণত address operator \texttt{\&} দিতে হয়, যেমন \texttt{\&n}।
\end{mistakebox}
